\chapter{A experiência da Caixa, 1907--1910}

\section{Do desenho legal à operação}

A aprovação da lei não garantia, por si, que o câmbio permaneceria próximo de 15 pence. A Caixa estava obrigada a receber ouro e emitir notas à taxa fixada, bem como a resgatar suas próprias notas pelo mesmo valor. Quando o mercado pressionava pela valorização do mil-réis, tornava-se vantajoso entregar ouro à instituição, pois ela pagava mais mil-réis por unidade de metal do que o mercado. A entrada ampliava reservas e circulação. Quando a pressão seguia o sentido contrário, portadores podiam resgatar notas e retirar ouro. A saída reduzia simultaneamente reservas e emissão.

A instituição operava apenas com suas próprias notas. O papel do Tesouro continuava inconversível, e o mercado cambial não desapareceu. A Caixa estabelecia uma referência e uma barreira à valorização, mas a defesa cotidiana da taxa também dependia das operações do Banco do Brasil. Oliveira e Silva mostram que o banco comprava e vendia divisas para manter o mercado próximo ao ponto legal, completando o mecanismo que a obrigação passiva de conversão não realizava sozinha \cite{oliveirasilva2001}.

Por esse motivo, a data do começo efetivo das operações e a divisão prática de tarefas entre Caixa, Tesouro e Banco do Brasil precisam ser reconstruídas com precisão. A lei foi sancionada em 6 de dezembro de 1906, mas o início administrativo, a emissão das primeiras notas e as intervenções do banco não devem ser fundidos num único marco.

\fonteaconferir{Localizar o regulamento, a data da abertura ao público, o primeiro balanço mensal e os primeiros registros de compra e venda do Banco do Brasil.}

Nos primeiros anos, o câmbio permaneceu próximo de 15 pence. Esse resultado combinou a regra legal, a ação do Banco do Brasil e a disponibilidade de recursos externos. Não demonstra que a Caixa isoladamente neutralizasse qualquer choque. A crise internacional de 1907 forneceu um primeiro teste dessa dependência.

\section{O pânico de 1907 e a valorização do café}

O pânico bancário iniciado nos Estados Unidos provocou recessão e retração do crédito internacional. Para o Brasil, o choque atingiu ao mesmo tempo exportações, financiamento e o esquema de valorização. O preço da borracha caiu, e a obtenção de recursos para carregar os estoques de café tornou-se mais difícil. O estado de São Paulo já havia comprado milhões de sacas, e o êxito da operação dependia de financiamento externo até que os estoques pudessem ser vendidos sem derrubar os preços \cite{fritsch1990,francolago2011}.

Em 1908, São Paulo obteve novo empréstimo de 15 milhões de libras, elevou a sobretaxa por saca e entregou aos credores maior controle sobre os estoques oferecidos em garantia. O governo federal concedeu garantia ao financiamento, reconhecendo a valorização como política de interesse nacional. A solução preservou o programa, mas distribuiu poder e risco entre o estado, a União, bancos e comerciantes estrangeiros. A política cafeeira que havia sido separada juridicamente da Caixa continuava ligada a ela pelos fluxos externos e pela taxa de câmbio.

A experiência de 1907 também relativiza uma narrativa linear de abundância entre 1906 e 1912. Houve uma interrupção séria, enfrentada por garantia pública, negociação com credores e recuperação posterior dos mercados. A estabilidade observada foi produzida sob tensão e com apoio institucional adicional.

\fonteaconferir{Confrontar o \textit{Retrospecto Commercial} de 1907 e 1908, os telegramas financeiros e a imprensa dos quatro jornais com a cronologia do empréstimo paulista e da garantia federal.}

\section{A taxa aduaneira e o significado prático da estabilidade}

Entre 1907 e 1909, o debate jornalístico não desapareceu. Mudou de objeto. Uma controvérsia importante dizia respeito ao valor utilizado para cobrar direitos aduaneiros. A libra valia no mercado cerca de 16 mil-réis, mas a Alfândega calculava determinados pagamentos por um padrão de 20 mil-réis. A diferença permitia discutir quem se beneficiava do padrão legal quando a prática econômica já se organizava em torno da taxa da Caixa.

Em julho de 1908, Barbosa Lima propôs cobrar os direitos segundo a taxa efetiva. O \textit{Correio da Manhã} apoiou a mudança e apresentou como escandalosa a existência de uma libra para bancos e comerciantes e outra para a Alfândega. O argumento deslocava o jornal de sua defesa de 1906 do padrão legal para a crítica ao descompasso entre lei e mercado. Também concedia que alterar o padrão não era, em si, desonestidade.

\fonteaconferir{Conferir a emenda Barbosa Lima, a sessão parlamentar de 14 de julho de 1908 e o artigo do \textit{Correio da Manhã}. Verificar se a formulação pertence a editorial próprio ou a comentário de sessão.}

O episódio mostra que o nível cambial produzia efeitos mesmo quando a taxa da Caixa não era formalmente alterada. Impostos, contratos em ouro, preços importados e operações do governo podiam conservar referências diferentes. O debate sobre estabilidade incluía, portanto, a unificação ou a coexistência de padrões de cálculo.

\section{A queda do horizonte de 27 pence}

Em agosto de 1907, o \textit{Correio da Manhã} publicou um artigo que registrava a perda de força do par de 27 e descrevia a Caixa como barreira ao retorno à antiga paridade. O texto chegou a sugerir que a taxa inicial deveria ter sido 16, e não 15. A mudança ocorreu menos de um ano depois da oposição severa do jornal ao projeto Campista.

\fonteaconferir{Conferir na página o artigo \textit{Melhorar pelo fogo}, de agosto de 1907, incluindo data, autoria e gênero. A extração tratada ainda não basta para sustentar citação literal.}

Há ao menos duas interpretações possíveis. A primeira é uma mudança doutrinária, na qual a experiência da Caixa teria convencido o jornal de que o retorno a 27 era inviável. A segunda é continuidade oposicionista e de interesses: o periódico alterava sua posição formal para criticar o governo e defender comércio ou lavoura diante da política vigente. As alternativas não são excludentes, e a trajetória completa precisa ser reconstruída peça por peça.

\interpretacaoprovisoria{O caso do \textit{Correio da Manhã} sugere que a direção numa escala abstrata pode mudar enquanto a função política do jornal conserva continuidades. A dissertação deverá distinguir revisão doutrinária, adaptação à nova realidade institucional e oposição ao governo de turno.}

O enfraquecimento de 27 não eliminou argumentos metalistas. O \textit{Correio Paulistano} continuou a justificar a Caixa como caminho de aproximação à circulação metálica. Quanto maior o ouro incorporado e as notas conversíveis emitidas, menor seria a proporção de papel inconversível no conjunto. A expansão da moeda conversível podia, nessa linguagem, realizar um fim ortodoxo.

\fonteaconferir{Reunir os artigos do \textit{Correio Paulistano} de 1907 a 1910 que relacionam crescimento dos depósitos, proporção de papel e circulação metálica.}

\textit{O Paiz} oferece outro cuidado metodológico. Em dezembro de 1908, chamou a Caixa de instituição insensata ao atacar a candidatura de David Campista. A hostilidade à instituição apareceu dentro de uma disputa sucessória. Antes de classificar a peça como doutrina monetária, é necessário separar crítica à política, ataque ao autor e estratégia de conjuntura.

\fonteaconferir{Conferir o editorial de \textit{O Paiz} de 21 de dezembro de 1908 e seu contexto sucessório.}

\section{Retorno dos capitais e expansão, 1908--1910}

Os mercados se recuperaram em 1908. A borracha voltou a contribuir para as receitas externas, empréstimos financiaram investimentos públicos e privados e capitais ingressaram no país. A retenção dos estoques ajudou a sustentar o café, cujos preços melhoraram a partir de 1909. Sob o mecanismo da Caixa, entradas de ouro se transformavam em aumento da circulação. A atividade econômica e a disponibilidade de crédito cresceram \cite{fritsch1990,francolago2011}.

O movimento evidencia o caráter pró-cíclico do arranjo. Boas exportações e entradas de capital elevavam reservas e moeda ao mesmo tempo. O aumento da renda estimulava importações e investimento, enquanto a fixação protegia exportadores de nova apreciação. O êxito imediato da estabilização não pode ser separado da conjuntura que alimentava os depósitos.

Em 1909 e nos primeiros meses de 1910, as condições do comércio e do financiamento externo foram excepcionalmente favoráveis. Segundo Oliveira e Silva, o Banco do Brasil encontrou dificuldade crescente para conter a valorização do mercado, e a taxa subiu acima de 15. Os depósitos na Caixa alcançaram 20 milhões de libras em maio de 1910, atingindo o limite que autorizava o Congresso a rever a paridade \cite{oliveirasilva2001}. A regra criada para uma possível valorização futura deixou de ser hipótese e se transformou em decisão imediata.

A chegada ao teto permite interpretar a reforma por duas perspectivas. Para os defensores da elevação, o papel inconversível e o mercado já valiam mais do que a nota conversível a 15, criando anomalia e pressão sobre o sistema. Para os opositores, a apreciação era justamente o movimento que a Caixa existia para deter. Elevar a taxa transferiria renda, reduziria a proteção à produção e abriria caminho a novas mudanças.

\section{A prática antes da lei em 1910}

Em abril de 1910, a mensagem presidencial e a exposição do ministro Leopoldo de Bulhões pediram mudança do regime. A proposta incluía elevação para 16, ampliação dos depósitos, possibilidade de novos aumentos e recomposição dos fundos monetários. No começo de maio, o Banco do Brasil passou a operar à taxa de 16 por ordem ministerial, antes de autorização legislativa. A medida criou uma situação em que prática bancária, Caixa e Alfândega podiam utilizar referências diferentes.

\fonteaconferir{Conferir a mensagem de Nilo Peçanha, a exposição de Bulhões e as ordens aplicadas pelo Banco do Brasil entre 4 e 6 de maio. Separar taxa de compra, taxa de venda e paridade legal.}

O episódio recolocou a questão da autoridade. Em 1906, a lei reservara ao Congresso a elevação depois de atingido o teto. Aplicar 16 por via administrativa permitia apresentar a mudança como resposta técnica ao mercado, mas também parecia antecipar uma decisão legislativa. A oposição pôde então combinar argumento econômico contra a valorização e argumento constitucional contra o Executivo.

\section{As posições em torno de 15 e 16 pence}

Cincinato Braga e Galeão Carvalhal figuraram entre os principais adversários da elevação. Sustentavam que 15 expressava melhor as condições econômicas e protegia produção e indústria. João Luiz Alves, relator da comissão de finanças do Senado, defendeu 16 e subordinou a divergência sobre o nível à importância maior da fixidez. Leopoldo de Bulhões apresentou a valorização como ajuste necessário diante do mercado e dos depósitos. A troca de atores em relação a 1906 é significativa: Campista perde centralidade, Bulhões retorna como protagonista e Calógeras ganha espaço no debate parlamentar.

Os jornais tornam visível a circulação desses argumentos, mas exigem distinção de voz. Em dezembro, \textit{O Paiz} publicou na \textit{Seção Livre} o discurso de Cincinato Braga contra aumento acima de 15. Dois dias depois publicou o parecer de João Luiz Alves, favorável a 16. Em seguida hospedou o discurso contrário de Galeão Carvalhal. Ao menos parte desse material ocupava espaço pago. A sequência documenta a composição pública do debate, não três mudanças editoriais em poucos dias.

\fonteaconferir{Conferir \textit{O Paiz} de 22, 24 e 29 de dezembro de 1910, identificar o estatuto de cada seção e comparar as reproduções dos mesmos discursos no \textit{Correio da Manhã} e na \textit{Gazeta de Notícias}.}

A voz própria de \textit{O Paiz} aparece com maior nitidez num editorial de maio. O texto elogiou Joaquim Murtinho, descreveu como paradoxo o papel inconversível valer mais do que a nota conversível e criticou a assimetria da Caixa, capaz de evitar depreciação além de um lado, mas incapaz de conter apreciação quando o limite era atingido. Essa apropriação é mais próxima da ortodoxia do que a simples publicação das seções livres.

\fonteaconferir{Conferir o editorial de \textit{O Paiz} de 14 de maio de 1910, especialmente a conclusão sobre a taxa e a relação com o projeto de Bulhões.}

O \textit{Correio da Manhã}, por sua vez, defendeu manter a Caixa no estado existente. Em 29 de maio, associou a reforma a ganhos de operadores cambiais e empresas endividadas em ouro, acusando Wenceslau Braz de sacrificar lavoura e indústria mineiras. A posição inverte a direção formal de 1906: o jornal que se opusera à fixação em 15 agora protegia essa taxa contra 16.

\fonteaconferir{Conferir o artigo \textit{Minas e o câmbio}, \textit{Correio da Manhã}, 29 de maio de 1910, e identificar a voz responsável.}

O \textit{Correio Paulistano} defendeu o aumento do ouro na circulação como avanço da metalicidade, argumento que atravessa a divisão simplificada entre expansão e ortodoxia. A posição da \textit{Gazeta de Notícias} em 1910 ainda precisa ser reconstruída em série. Seu editorial retrospectivo de janeiro de 1912 afirmou que o jornal havia auxiliado a Caixa e se associado aos protestos contra a ação de Bulhões, mas essa memória posterior não substitui a leitura contemporânea da reforma.

\lacuna{Completar a janela de abril a dezembro de 1910 na \textit{Gazeta de Notícias} e confrontar o resultado com a autodescrição publicada em 15 de janeiro de 1912.}

\section{Fundos, limite e novo desenho legal}

O debate de 1910 não tratou apenas da taxa. Parlamentares questionaram a existência e o destino do fundo de garantia. Requerimentos de informação foram recusados, e o saldo escriturado não parecia corresponder a recursos imediatamente disponíveis. A disputa retomava um problema de 1906: se os fundos eram instrumentos de valorização futura, lastro, reserva para intervenção ou recursos consumidos na sustentação do câmbio.

\fonteaconferir{Localizar os requerimentos de informação, as recusas nas duas casas e a composição do saldo de 8.069.093 libras mencionada nas peças tratadas.}

O Decreto n. 2.357, de 31 de dezembro, elevou a taxa da Caixa para 16 pence, ampliou o limite de emissão para 900 mil contos, correspondente a 60 milhões de libras, e restaurou os fundos de resgate e garantia \cite{brasil1910taxa16}. O novo teto triplicou a cifra em mil-réis da lei original. A reforma preservou a conversibilidade e respondeu à valorização por meio de nova expansão potencial, em vez de encerrar a Caixa quando o primeiro limite foi alcançado.

O compromisso de 1910 diferiu do de 1906. A taxa maior atendia aos que julgavam anômala a diferença entre mercado e paridade. O limite ampliado permitia que o mecanismo continuasse absorvendo ouro. A restauração dos fundos procurava manter um horizonte de garantia e resgate. Ao mesmo tempo, o debate deixou exposta a fragilidade da informação sobre reservas e a tensão entre decisão legislativa e prática administrativa.

\interpretacaoprovisoria{A reforma marcou o auge institucional da Caixa, mas também revelou que a fixidez dependia de ajustes políticos. Quando o teto original se tornou vinculante, o regime não produziu uma solução automática: governo e Congresso escolheram nova taxa, novo limite e nova composição de fundos. A estabilidade era uma regra administrada por decisões sucessivas.}

O período seguinte colocaria uma pergunta diferente. Até 1910, a pressão dominante vinha da entrada de ouro e da valorização. A partir de 1912 e 1913, exportações, empréstimos e reservas se moveriam no sentido contrário. O problema deixaria de ser quanto expandir para absorver abundância e passaria a ser quanto contrair, que compromissos preservar e quem suportaria a defesa da paridade.
